% =====================================================================
%  08_conclusiones.tex — Conclusiones (preliminares en esta versión)
% =====================================================================
\chapter{Conclusiones}
\label{cap:conclusiones}

\begin{nota}[title={Conclusiones preliminares}]
Estas conclusiones corresponden al Prototipo~1 y se actualizarán al cerrar
cada prototipo.
\end{nota}

El Prototipo~1 de \saru{} demuestra que es viable registrar de forma
automática las entradas y salidas de unidades en la puerta del Patio CCA
con procesamiento en el borde. Una sola Jetson ejecuta tres modelos en
cascada, decide los cruces, registra los movimientos, publica un tablero y
graba video continuo con un consumo de entre 11 y 19~W.

La metodología de prototipado evolutivo resultó adecuada: varios de los
casos más importantes, como las unidades detenidas sobre la línea, varias
unidades a la vez o los números leídos con dígitos de menos, solo
aparecieron al probar con video real, y se corrigieron sin detener la
operación.

Los principales retos para los siguientes prototipos son tres: garantizar
la entrega de los avisos desde la red del patio, mejorar la lectura en
vistas no laterales y medir formalmente la exactitud del sistema. A ellos
se suma un reto del proyecto: unificar las versiones del código y mantener
esta documentación al día con cada cambio.
